% Part: many-valued-logic
% Chapter: three-valued-logics
% Section: multiple-designation

\documentclass[../../../include/open-logic-section]{subfiles}

\begin{document}

\olfileid{mvl}{thr}{mul}

\olsection{ممتاز دانستن ارزشی جز~$\True$}

همهٔ منطق‌هایی که تاکنون دیده‌ایم مجموعهٔ ارزش‌های صدق ممتاز
$V^+ = \{\True\}$ را داشته‌اند؛ یعنی چیزی صادق شمرده می‌شود اگر و تنها
اگر ارزش صدقش~$\True$ باشد. اما می‌توان چیزی را صادق شمرد اگر فقط
$\False$ نباشد. در آن صورت، برای مثال با مقررکردن
$V^+ = \{\True, \Undef\}$ در ماتریس، منطقی به دست می‌آید.

\begin{defn}
\emph{منطق پارادوکس}~$\LogLP$ با ماتریس زیر تعریف می‌شود:
\begin{enumerate}
  \item زبان گزاره‌ای استاندارد~$\Lang L_0$ با~$\lnot$، $\land$،
  $\lor$ و~$\lif$.
  \item مجموعهٔ ارزش‌های صدق~$V = \{\True, \Undef, \False\}$.
  \item $\True$ و~$\Undef$ ممتازند؛ یعنی~$V^+ = \{\True, \Undef\}$.
  \item توابع صدق همان توابع منطق قوی کلینی‌اند.
\end{enumerate}
\end{defn}

\begin{defn}
\emph{منطق بی‌معنایی} هالدن~$\LogHal$ با ماتریس زیر تعریف می‌شود:
\begin{enumerate}
  \item زبان گزاره‌ای استاندارد~$\Lang L_0$ با~$\lnot$، $\land$،
  $\lor$، $\lif$ و رابط یک‌موضعی~$+$.
  \item مجموعهٔ ارزش‌های صدق~$V = \{\True, \Undef, \False\}$.
  \item $\True$ و~$\Undef$ ممتازند؛ یعنی~$V^+ = \{\True, \Undef\}$.
  \item توابع صدق همان توابع منطق ضعیف کلینی‌اند، به‌اضافهٔ عملگر
  «بی‌معناست»:
  \begin{center}
    \begin{tabular}{c|c} 
    $\tf{+}$ & \\ 
    \hline  
    $\True$ & $\False$ \\ 
    $\Undef$ & $\True$ \\
    $\False$ & $\False$ 
    \end{tabular}
  \end{center}
\end{enumerate}
\end{defn}

برخلاف منطق‌های کلینی که جدول‌های صدق مشترکی با آن‌ها دارند، این منطق‌ها
\emph{دارای} همان‌گویی‌اند.

\begin{prop}\ollabel{prop:LP-taut-CL} همان‌گویی‌های~$\LogLP$ همان
  همان‌گویی‌های منطق گزاره‌ای کلاسیک‌اند.
\end{prop}

\begin{proof}
  بنا بر~\olref[syn][sub]{prop:mvl-cl}، اگر~$\Entails[\LogLP] !A$،
  آنگاه~$\Entails[\LogCL] !A$. برای نشان‌دادن جهت عکس، ثابت می‌کنیم اگر
  !!a{valuation} چون
  $\pAssign v\colon \PVar \to \{\False, \True, \Undef\}$ وجود داشته
  باشد که~$\pValue v[\LogKs](!A) = \False$، آنگاه~!!a{valuation} چون
  $\pAssign {v'}\colon \PVar \to \{\False, \True\}$ وجود دارد که
  $\pValue {v'}[\LogCL](!A) = \False$. این ادعا نتیجه را برای~$\LogLP$
  ثابت می‌کند، زیرا~$\LogKs$ و~$\LogLP$ توابع صدق مشخصهٔ یکسان دارند و
  $\False$ تنها ارزش صدق ناممتاز~$\LogLP$ است؛ همین تنها تفاوت~$\LogLP$
  و~$\LogKs$ است. پس اگر~$\Entails/[\LogLP] !A$، برای~!!{valuation}ای
  چون~$\pAssign v$ داریم
  $\pValue v[\LogLP](!A) = \pValue v[\LogKs](!A) = \False$. بنا بر
  ادعایی که ثابت می‌کنیم، $\pValue{v'}[\LogCL](!A) = \False$؛ یعنی
  $\Entails/[\LogCL] !A$.
  
  برای اثبات ادعا، نخست~$\pAssign {v'}$ را چنین تعریف می‌کنیم:
  \[
    \pAssign {v'}(p) = 
    \begin{cases}
    \True & \text{اگر } \pAssign {v}(p) \in \{\True, \Undef\}\\
    \False & \text{در غیر این صورت}
  \end{cases}
  \]
  اکنون با استقرا بر~$!A$ نشان می‌دهیم که (الف) اگر
  $\pValue v[\LogKs](!A) = \False$، آنگاه
  $\pValue {v'}[\LogCL](!A) = \False$؛ و (ب) اگر
  $\pValue v[\LogKs](!A) = \True$، آنگاه
  $\pValue {v'}[\LogCL](!A) = \True$.
  \begin{enumerate}
    \item پایهٔ استقرا: $!A \ident p$. بنا بر
    \olref[syn][val]{defn:pValue}، اگر
    $\pValue v[\LogKs](!A) = \False$، آنگاه
    $\pAssign v(p) = \False$ و از تعریف~$v'$ داریم
    $\pValue {v'}[\LogCL](!A) = \False$. همچنین اگر
    $\pValue v[\LogKs](!A) = \True$، آنگاه
    $\pAssign v(p) = \True$ و از تعریف~$v'$ داریم
    $\pValue {v'}[\LogCL](!A) = \True$. پس هر دو بند (الف) و (ب)
    برقرارند.
    
    برای گام استقرا، حالت‌های زیر را در نظر می‌گیریم:
    \item $!A \ident \lnot !B$.
    \begin{enumerate}
      \item فرض کنید~$\pValue v[\LogKs](\lnot!B) = \False$. بنا بر
      تعریف~$\tf{\lnot}[\LogKs]$،
      $\pValue v[\LogKs](!B) = \True$. بنا بر فرض استقرا، بند~(ب)،
      $\pValue {v'}[\LogCL](!B) = \True$؛ پس
      $\pValue {v'}[\LogCL](\lnot !B) = \False$.
      \item فرض کنید~$\pValue v[\LogKs](\lnot!B) = \True$. بنا بر
      تعریف~$\tf{\lnot}[\LogKs]$،
      $\pValue v[\LogKs](!B) = \False$. بنا بر فرض استقرا، بند~(الف)،
      $\pValue {v'}[\LogCL](!B) = \False$؛ پس
      $\pValue {v'}[\LogCL](\lnot !B) = \True$.
    \end{enumerate}

    \item $!A \ident (!B \land !C)$.
    \begin{enumerate}
      \item فرض کنید~$\pValue v[\LogKs](!B \land !C) = \False$. بنا بر
      تعریف~$\tf{\land}[\LogKs]$، یا
      $\pValue v[\LogKs](!B) = \False$، یا
      $\pValue v[\LogKs](!C) = \False$. بنا بر فرض استقرا، بند~(الف)،
      یا~$\pValue {v'}[\LogCL](!B) = \False$، یا
      $\pValue {v'}[\LogCL](!C) = \False$؛ پس
      $\pValue {v'}[\LogCL](!B \land !C) = \False$.
      \item فرض کنید~$\pValue v[\LogKs](!B \land !C) = \True$. بنا بر
      تعریف~$\tf{\land}[\LogKs]$،
      $\pValue v[\LogKs](!B) = \True$ و
      $\pValue v[\LogKs](!C) = \True$. بنا بر فرض استقرا، بند~(ب)،
      $\pValue {v'}[\LogCL](!B) = \True$ و
      $\pValue {v'}[\LogCL](!C) = \True$؛ پس
      $\pValue {v'}[\LogCL](!B \land !C) = \True$.
    \end{enumerate}
  \end{enumerate}

  دو حالت دیگر مشابه‌اند و به‌عنوان تمرین واگذار می‌شوند. راه دیگر این
  است که برهان بالا را اثبات نتیجه برای همهٔ~!!{formula}هایی بدانیم که
  فقط شامل~$\lnot$ و~$\land$ هستند. سپس می‌توان از این واقعیت استفاده
  کرد که هم در~$\LogKs$ و هم در~$\LogCL$، برای هر~$\pAssign v$،
  $\pValue v(!B \lor !C) = \pValue v(\lnot(\lnot!B \land \lnot !C))$
  و~$\pValue v(!B \lif !C) = \pValue v(\lnot(!B \land \lnot !C))$.
\end{proof}

\begin{prob}
  برهان~\olref[mvl][thr][mul]{prop:LP-taut-CL} را کامل کنید؛ یعنی
  بندهای~(الف) و~(ب) را برای حالت‌های~$!A \ident (!B \lor !C)$ و
  $!A \ident (!B \lif !C)$ ثابت کنید.
\end{prob}

\begin{prob}
ثابت کنید هر همان‌گویی کلاسیک در~$\LogHal$ همان‌گویی است.
\end{prob}

با اینکه همان‌گویی‌های این منطق‌ها همان همان‌گویی‌های منطق کلاسیک‌اند،
روابط نتیجهٔ آن‌ها متفاوت است. برای نمونه، $\LogLP$
\emph{فرا‌سازگار} است، زیرا~$\lnot p, p \Entails/ q$؛ پس اصل انفجار
$\lnot !A, !A \Entails !B$ به‌طور کلی برقرار نیست. (این اصل برای برخی
حالت‌های~$!A$ و~$!B$ برقرار است؛ مثلاً اگر~$!B$ همان‌گویی باشد.)

\begin{prob}
  کدام‌یک از روابط زیر در (الف)~$\LogLP$ و (ب)~$\LogHal$ برقرارند؟
  برای هریک جدول صدق بدهید.
  \begin{enumerate}
    \item $p, p \lif q \Entails q$
    \item $\lnot q, p \lif q \Entails \lnot p$
    \item $p \lor q, \lnot p \Entails q$
    \item $\lnot p, p \Entails q$
    \item $p \Entails p \lor q$
    \item $p \lif q, q\lif r \Entails p \lif r$
  \end{enumerate}
\end{prob}

اگر~$\Undef$ را در~$\LogLuk[3]$ ممتاز کنیم چه می‌شود؟

\begin{defn}
  \emph{R-Mingle سه‌ارزشی}~$\LogRM[3]$ با ماتریس زیر تعریف می‌شود:
  \begin{enumerate}
    \item زبان گزاره‌ای استاندارد~$\Lang L_0$ با~$\lfalse$، $\lnot$،
    $\land$، $\lor$ و~$\lif$.
    \item مجموعهٔ ارزش‌های صدق~$V = \{\True, \Undef, \False\}$.
    \item $\True$ و~$\Undef$ ممتازند؛ یعنی~$V^+ = \{\True, \Undef\}$.
    \item توابع صدق همان توابع منطق \L ukasiewicz، یعنی~$\LogLuk[3]$،
    هستند.
  \end{enumerate}
\end{defn}

\begin{prob}
  کدام‌یک از روابط زیر در~$\LogRM[3]$ برقرارند؟
  \begin{enumerate}
    \item $p, p \lif q \Entails q$
    \item $p \lor q, \lnot p \Entails q$
    \item $\lnot p, p \Entails q$
    \item $p \Entails p \lor q$
  \end{enumerate}
\end{prob}

گاهی جدول‌های صدق متفاوت فقط با تغییر ارزش‌های ممتاز، یک منطق واحد،
یعنی یک رابطهٔ استلزام واحد، تولید می‌کنند. برای مثال، اگر در منطق گودل
به‌جای~$\{\True\}$، مجموعهٔ~$V^+ = \{\True, \Undef\}$ را بگیریم، چنین
می‌شود.

\begin{prop}\ollabel{prop:gl-udes}
  ماتریسی با~$V = \{\False, \Undef,\True\}$،
  $V^+=\{\True, \Undef\}$ و توابع صدق منطق سه‌ارزشی گودل، منطق کلاسیک
  را تعریف می‌کند.
\end{prop}

\begin{proof}
  تمرین.
\end{proof}

\begin{prob}
  \olref[mvl][thr][mul]{prop:gl-udes} را چنین ثابت کنید: برای منطق
  $\Log L$ که درست مانند منطق گودل، اما با
  $V^+=\{\True,\Undef\}$، تعریف شده است، نشان دهید اگر
  $\Gamma \Entails/[\Log L] !B$، آنگاه~$\Gamma \Entails/[\LogCL] !B$.
  از اندیشه‌های~\olref[mvl][thr][mul]{prop:LP-taut-CL} استفاده کنید، با
  این تفاوت که به‌جای اثبات خواص~(الف) و~(ب)، نشان دهید
  $\pValue v[\LogGod](!A) = \False$ اگر و تنها اگر
  $\pValue {v'}[\LogCL](!A) = \False$؛ و در نتیجه
  $\pValue v[\LogGod](!A) \in \{\True,\Undef\}$ اگر و تنها اگر
  $\pValue {v'}[\LogCL](!A) = \True$. توضیح دهید چرا این نتیجه گزاره
  را ثابت می‌کند.
\end{prob}

\end{document}
